\section{I/O and Streams}

\subsection{Section Overview}
So far, we've used \texttt{std::cin} and \texttt{std::cout} for basic console input and output. These are part of the larger C++ I/O stream library. In this section, we'll explore this library in more detail, learning how to format our output with stream manipulators, how to read from and write to files, and how to use string streams for in-memory formatting.

\subsection{Files, Streams and I/O}
The C++ I/O library is based on the concept of a \textbf{stream}, which is a sequence of bytes flowing from a source to a destination.
\begin{itemize}
    \item An \textbf{input stream} (\texttt{istream}) reads bytes from a source (like the keyboard or a file).
    \item An \textbf{output stream} (\texttt{ostream}) writes bytes to a destination (like the console or a file).
\end{itemize}
The \texttt{<iostream>} header gives us \texttt{cin} and \texttt{cout}. The \texttt{<fstream>} header gives us file streams (\texttt{ifstream} and \texttt{ofstream}). The \texttt{<sstream>} header gives us string streams (\texttt{stringstream}).

\subsection{Stream Manipulators --- Boolean}
Stream manipulators are used to control the formatting of I/O. For booleans, you can display them as \texttt{1} and \texttt{0} or \texttt{true} and \texttt{false}.
\begin{lstlisting}
#include <iostream>
#include <iomanip> // required for manipulators

int main() {
	std::cout << "Default: " << true << " " << false << std::endl;
	std::cout << std::boolalpha; // display as true/false
	std::cout << "boolalpha: " << true << " " << false << std::endl;
	std::cout << std::noboolalpha; // back to 1/0
	std::cout << "noboolalpha: " << true << " " << false << std::endl;
}
\end{lstlisting}

\subsection{Stream Manipulators --- Integer}
You can control the base (decimal, hexadecimal, octal) and show the base prefix.
\begin{lstlisting}
int num = 255;
std::cout << "Decimal: " << std::dec << num << std::endl;
std::cout << "Hexadecimal: " << std::hex << num << std::endl;
std::cout << "Octal: " << std::oct << num << std::endl;

// show the base prefix
std::cout << std::showbase;
std::cout << "Hex with base: " << std::hex << num << std::endl; // 0xff
\end{lstlisting}

\subsection{Stream Manipulators --- Floating Point}
You can control the precision, notation (scientific, fixed), and whether to show the decimal point.
\begin{lstlisting}
double num = 1234.5678;
std::cout << std::setprecision(5) << num << std::endl; // 1234.6
std::cout << std::fixed << std::setprecision(2) << num << std::endl; // 1234.57
std::cout << std::scientific << num << std::endl; // 1.23e+03
\end{lstlisting}

\subsection{Stream Manipulators --- Align and Fill}
\texttt{setw(width)} sets the field width for the next output operation. \texttt{left}, \texttt{right}, and \texttt{internal} control alignment. \texttt{setfill(char)} sets the fill character.
\begin{lstlisting}
double num = 12.34;
std::cout << std::setw(10) << num << std::endl; // "     12.34"
std::cout << std::setw(10) << std::left << num << std::endl; // "12.34     "
std::cout << std::setw(10) << std::setfill('-') << std::left << num << std::endl; // "12.34-----"
\end{lstlisting}

\subsection{Reading from a Text File}
To read from a file, you use \texttt{std::ifstream}. You must include the \texttt{<fstream>} header.
\begin{lstlisting}
#include <fstream>
#include <iostream>

int main() {
	std::ifstream in_file("my_data.txt");
	if (!in_file) {
		std::cerr << "Error opening file" << std::endl;
		return 1;
	}

	int num;
	std::string line;
	in_file >> num >> line; // read an int and a word

	std::cout << "Read: " << num << " " << line << std::endl;

	in_file.close();
	return 0;
}
\end{lstlisting}

\subsection{Coding Exercise 36: Reading from a Text File}
Given a file \texttt{data.txt} that contains a list of numbers, one per line, write a program that reads all the numbers and prints their sum.

\textbf{File \texttt{data.txt}:}
\begin{lstlisting}[numbers=none]
10
20
30
\end{lstlisting}
\textbf{Solution:}
\begin{lstlisting}
#include <fstream>
#include <iostream>

int main() {
	std::ifstream in_file("data.txt");
	int sum = 0, num;
	while (in_file >> num) {
		sum += num;
	}
	std::cout << "Sum: " << sum << std::endl;
	in_file.close();
	return 0;
}
\end{lstlisting}

\subsection{Writing to a Text File}
To write to a file, you use \texttt{std::ofstream}.
\begin{lstlisting}
#include <fstream>

int main() {
	std::ofstream out_file("output.txt"); // creates the file
	// std::ofstream out_file("output.txt", std::ios::app); // append mode

	if (!out_file) {
		std::cerr << "Error creating file" << std::endl;
		return 1;
	}

	out_file << "Hello, file!" << std::endl;
	out_file << 123 << std::endl;

	out_file.close();
	return 0;
}
\end{lstlisting}

\subsection{Coding Exercise 37: Writing to a Text File}
Write a program that writes your name and age to a file named \texttt{info.txt}.
\begin{lstlisting}
#include <fstream>
#include <string>

int main() {
	std::ofstream out_file("info.txt");
	std::string name = "John Doe";
	int age = 30;
	out_file << "Name: " << name << std::endl;
	out_file << "Age: " << age << std::endl;
	out_file.close();
	return 0;
}
\end{lstlisting}

\subsection{Using stringstreams}
\texttt{stringstream} allows you to use stream insertion and extraction operators on strings in memory. This is very useful for formatting strings or for parsing data from a string. You must include the \texttt{<sstream>} header.
\begin{lstlisting}
#include <sstream>

int main() {
	std::stringstream ss;
	int age = 30;
	std::string name = "Alice";
	ss << "Name: " << name << ", Age: " << age;

	std::string formatted_string = ss.str();
	std::cout << formatted_string << std::endl;

	// Using stringstream to parse
	std::stringstream parser("100 200.5");
	int a;
	double b;
	parser >> a >> b;
	std::cout << a << " " << b << std::endl;
}
\end{lstlisting}

\subsection{Section Challenge 1}
Write a program to format and display a table of data. Use stream manipulators to align the columns and set the precision of floating-point numbers.

\subsection{Section Challenge 1 --- Solution}
\begin{lstlisting}
#include <iostream>
#include <iomanip>

int main() {
	std::cout << std::left;
	std::cout << std::setw(20) << "Name" << std::setw(10) << "Score" << std::endl;
	std::cout << std::setw(30) << std::setfill('-') << "" << std::endl;
	std::cout << std::setfill(' ');

	std::cout << std::setw(20) << "Alice" << std::setw(10) << std::right << 95.5 << std::endl;
	std::cout << std::setw(20) << std::left << "Bob" << std::setw(10) << std::right << 88.0 << std::endl;

	return 0;
}
\end{lstlisting}

\subsection{Section Challenge 2}
Write a program that reads a text file, where each line contains a student's name and their score (e.g., ``Alice 95''). The program should calculate the average score and print a report.

\subsection{Section Challenge 2 --- Solution}
\textbf{File \texttt{scores.txt}:}
\begin{lstlisting}[numbers=none]
Alice 95
Bob 88
Charlie 92
\end{lstlisting}
\textbf{Solution:}
\begin{lstlisting}
#include <iostream>
#include <fstream>
#include <string>
#include <vector>

int main() {
	std::ifstream in_file("scores.txt");
	std::string name;
	int score;
	double total = 0;
	int count = 0;

	while (in_file >> name >> score) {
		std::cout << "Read: " << name << " " << score << std::endl;
		total += score;
		count++;
	}

	if (count > 0) {
		std::cout << "Average score: " << total / count << std::endl;
	}

	in_file.close();
	return 0;
}
\end{lstlisting}

\subsection{Section Challenge 3}
Write a program that reads a file and creates a copy of it, but with line numbers added at the beginning of each line.

\subsection{Section Challenge 3 --- Solution}
\begin{lstlisting}
#include <iostream>
#include <fstream>
#include <string>

int main() {
	std::ifstream in_file("original.txt");
	std::ofstream out_file("numbered_copy.txt");
	std::string line;
	int line_number = 1;

	if (!in_file || !out_file) {
		std::cerr << "Error with files." << std::endl;
		return 1;
	}

	while (std::getline(in_file, line)) {
		out_file << std::setw(5) << std::left << line_number++ << line << std::endl;
	}

	std::cout << "File copied with line numbers." << std::endl;
	in_file.close();
	out_file.close();
	return 0;
}
\end{lstlisting}

\subsection{Section Challenge 4}
Write a program that reads a file word by word, and finds and prints the number of times a specific word (provided by the user) appears in the file.

\subsection{Section Challenge 4 --- Solution}
\begin{lstlisting}
#include <iostream>
#include <fstream>
#include <string>

int main() {
	std::ifstream in_file("story.txt");
	std::string word_to_find;
	std::cout << "Enter the word to find: ";
	std::cin >> word_to_find;

	int count = 0;
	std::string current_word;
	while (in_file >> current_word) {
		if (current_word == word_to_find) {
			count++;
		}
	}

	std::cout << "The word '" << word_to_find << "' was found " << count << " times." << std::endl;
	in_file.close();
	return 0;
}
\end{lstlisting}

\subsection{Quiz 16: Section 19 Quiz}
\begin{enumerate}
    \item Which header file is needed for file I/O operations like \texttt{ifstream} and \texttt{ofstream}?
    \begin{enumerate}
        \item[a)] \texttt{<iostream>}
        \item[b)] \texttt{<stream>}
        \item[c)] \texttt{<fstream>}
        \item[d)] \texttt{<sstream>}
    \end{enumerate}

    \item What is the purpose of \texttt{std::boolalpha}?
    \begin{enumerate}
        \item[a)] To format integers as hexadecimal.
        \item[b)] To set the field width for output.
        \item[c)] To format boolean values as ``true'' and ``false''.
        \item[d)] To align output to the right.
    \end{enumerate}

    \item What is \texttt{stringstream} useful for?
    \begin{enumerate}
        \item[a)] Reading from files.
        \item[b)] Writing to files.
        \item[c)] Performing formatted I/O operations on strings in memory.
        \item[d)] Handling network sockets.
    \end{enumerate}
\end{enumerate}

\textbf{Answers:} 1. (c), 2. (c), 3. (c)
